\documentclass[sigconf, nonacm]{acmart}


\usepackage{draftwatermark}
\SetWatermarkText{AI-Generated}
\SetWatermarkScale{0.5}
\SetWatermarkColor[gray]{0.6}
\AtBeginDocument{\providecommand\BibTeX{{Bib\TeX}}}

\copyrightyear{2025}
\acmYear{2025}
\setcopyright{none}
\acmConference[arXiv preprint]{arXiv preprint}{2025}{}
\acmBooktitle{}
\acmDOI{}
\acmISBN{}
\settopmatter{printacmref=false}         
\renewcommand\footnotetextcopyrightpermission[1]{} 

\usepackage{multirow}
\usepackage{enumitem}
\usepackage[skip=2pt]{caption}
\usepackage{bm}
\graphicspath{{figures/}}


\begin{document}
\title{FoodLogAthl-218: Constructing a Real-World Food Image Dataset\\ Using Dietary Management Applications}


\author{Anonymous Author}

\begin{abstract}
Food recognition models for dietary logging are usually evaluated on web-crawled images that are attractive, class-balanced, independent and show one dish, whereas photos taken for personal meal records are casual, long-tailed, temporally ordered and often contain several dishes. We introduce FoodLogAthl-218, a benchmark built bottom-up from meal photos submitted to the FoodLog Athl dietary management application: labels follow what users logged and are consolidated afterwards. Automated filtering, CLIP-based outlier detection, manual review with removal of personal information, and GPT-4o-assisted label merging with human adjudication turn 11,460 raw images into 6,925 images, 14,349 dish bounding boxes and 218 classes from 513 users, with anonymized user identifiers and timestamps. Three tasks are defined: closed-set dish classification, continual user-aware adaptation along the time axis, and menu-aware classification in multi-dish images. Ten off-the-shelf vision-language models reach at best 42.4\% top-1 recall (GPT-4o) and 60.0\% top-5 recall (CLIP ViT-L/14), and feature statistics show larger intra-class spread and class overlap than UEC-Food 256. A Qwen2.5-VL-3B model fine-tuned along the timeline improves from roughly 0.25 to 0.5 accuracy while the zero-shot model does not, and a user-specific adapter gains about 50 points for one heavy logger. Menu descriptions raise F1 from 0.159 to 0.230 but lower top-1 recall from 0.366 to 0.324. The data come from one application and inter-annotator agreement was not measured, so the benchmark is a measurement tool for realistic food recognition, not a comprehensive one.
\end{abstract}

\maketitle
\section{Introduction}
Food recognition is a building block of image-based dietary logging, where a user photographs a meal and the system proposes dishes to record. Progress has been measured mostly on benchmarks assembled by crawling the web for images of predefined classes~\cite{bossard2014food,kawano2015automatic,min2023large}. Such data differ from what a logging application receives in four ways (Figure~\ref{fig:overview}): (i) the label set is fixed first and images are collected afterwards, whereas in a logging application the labels follow what people eat; (ii) class frequencies are balanced or curated, whereas real consumption is long-tailed; (iii) web images are made to be shared and look appealing, whereas logging photos are casual and unfiltered; and (iv) samples are independent, whereas a personal log is a time series of one person's habits, usually with several dishes per photo. A model that scores well on a web benchmark need not behave equally well in deployment.

\begin{figure*}[!tb]
  \centering
  \includegraphics[width=\textwidth]{distribution_5.jpg}
  \caption{Overview: traditional web-collected food image datasets (top) compared with FoodLogAthl-218 (bottom) in terms of data source, labelling order, class balance, aesthetic bias, temporal continuity and number of dishes per photo.}
  \label{fig:overview}
\end{figure*}

We present FoodLogAthl-218, a benchmark built from the records of the FoodLog Athl application~\cite{nakamoto2022foodlog,aizawa2015foodlog}. Our design principles are \emph{realism} (natural imbalance, casual photographs, multi-dish meals, temporal order), \emph{label quality after the fact} (users' free choices are cleaned and consolidated by outlier detection, language-model label grouping and human adjudication) and \emph{diagnostic value} (three tasks separating closed-set recognition, adaptation over time and use of meal context). The final data contain 6,925 images, 14,349 dish bounding boxes and 218 classes. We evaluate contrastive vision-language models and large multimodal models (LMMs) and ask how hard closed-set classification on real logs is, how much models gain from adaptation along the timeline or to one user, and whether the rest of the menu helps classify a single dish.

Our contributions are: (1) a real-world food image dataset with user, time and meal metadata, natural imbalance and multi-dish images; (2) a mixed automated--human curation pipeline with per-step statistics; (3) three benchmark tasks with a protocol and metrics; (4) a validation of diversity and difficulty against UEC-Food 256; and (5) baseline results showing low closed-set accuracy, a large benefit of temporal and personal adaptation in our setting, and a precision--recall trade-off when menu context is used.

\section{Related Works}
\textbf{Food image datasets.} Early food datasets covered one to a few hundred classes of mostly Western or fast food~\cite{chen2009pfid,bossard2014food}; later ones broadened cuisine and scale, e.g., UEC-Food 256~\cite{kawano2015automatic}, UPMC Food-101~\cite{wang2015recipe}, FoodAI-756~\cite{sahoo2019foodai}, Food2K~\cite{min2023large} and Food-500 Cap~\cite{ma2023food}, or target vegetables and fruits~\cite{hou2017vegfru}, participants' own photos~\cite{farinella2016retrieval} and recipes~\cite{marin2019learning,harashima2017cookpad}. Table~\ref{tab:datasets} contrasts these with our dataset: almost all are web-collected with predefined classes. Other work adds segmentation~\cite{wu2021large}, ingredients~\cite{chen2020study} or nutrition~\cite{thames2021nutrition5k}. FoodLogAthl-218 is smaller than the large web datasets; its distinguishing property is that images and labels originate from dietary logging itself.

\textbf{Personal food logs and personalization.} Food-logging applications~\cite{aizawa2015foodlog,nakamoto2022foodlog,asken,myfitnesspal} and mobile food diaries~\cite{meyers2015im2calories} motivated work on personal dietary tendencies~\cite{Aizawa2013FoodBalance} and personalized classifiers~\cite{flg469}. These works use proprietary logs; a public benchmark with user and time information for standardized evaluation of personalization has been missing.

\textbf{Vision-language and multimodal models for food.} CLIP-style models~\cite{CLIP,sun2023eva} allow zero-shot classification against label text, and LMMs~\cite{liu2024visual,liu2024improved,bai2023qwen} can describe a dish in free text. Food-specific LMM work and benchmarks include FoodLMM~\cite{yin2023foodlmm}, FoodMLLM-JP~\cite{Imajuku2025FoodMLLM} and WorldCuisines~\cite{winata2024worldcuisines}, which targets multilingual visual question answering. We use these model families as baselines through lmms-eval~\cite{lmms_eval2024,zhang2024lmmsevalrealitycheckevaluation}.

\section{FoodLogAthl-218 Dataset}
\subsection{Data Source}
FoodLog Athl is a dietary management application in which users photograph meals; the application proposes dish names from a database of about 140K items, detects the dishes with bounding boxes, and lets users correct the labels and record portion sizes~\cite{nakamoto2022foodlog}. We took the logs from May 2023 to October 2024: 11,460 images with 61,414 bounding boxes and 5,889 raw labels from 632 users, each with anonymized user ID, date and time, dish names, portion size and the initial detections. Because users may or may not correct labels, the raw data contain three kinds of errors: incorrect labels, incorrect bounding boxes (a box placed on the wrong dish) and inappropriate images (e.g., an empty plate labelled as milk, or a food package label labelled as a donut). The image is the atomic unit for filtering, so that multi-dish meal composition is preserved.

\subsection{Curation Pipeline}
Table~\ref{tab:filtering} gives the statistics of each step. (1--2) Samples without image or bounding box are removed, leaving 11,285 images, 26,065 bounding boxes and 3,071 classes. (3) Classes with fewer than 10 samples are dropped (9,121 images, 21,944 bounding boxes, 1,180 classes). (4) CLIP ViT-B/16~\cite{CLIP} features are used to flag as outliers the samples more than 1.5 standard deviations from their class centroid, which leaves 7,537 images and 17,513 bounding boxes. (5) Reviewers manually fix or remove invalid boxes and labels; an image is removed entirely (all-or-none) when any of its boxes is invalid, so that meal composition is not distorted, resulting in 6,925 images, 14,349 bounding boxes and 513 users. (6) GPT-4o groups synonymous labels, and manual refinement merges them into the final 218 classes. Personal information was removed in two manual stages, and the data collection has ethics approval and informed consent from users. Only 60\% of the raw images and 23\% of the raw bounding boxes survive.

\input{tables/table_zn}

\subsection{Dataset Characteristics}
The class distribution is long-tailed (Figure~\ref{fig:hist}): the number of samples per class ranges from the filtering threshold of 10 to well over a thousand, with many classes between 10 and 100. Table~\ref{tab:datasets} compares FoodLogAthl-218 with existing food classification datasets. Note that the 14,349 samples of our dataset are dish bounding boxes from 6,925 images, whereas the other datasets are listed with the sample counts reported in their papers.

\begin{figure}[!tb]
  \centering
  \includegraphics[width=0.8\columnwidth]{histogram_linear_pdf.jpg}
  \caption{Number of dish types (classes) per number of images in FoodLogAthl-218 (logarithmic horizontal axis), with illustrative examples marked by arrows.}
  \label{fig:hist}
\end{figure}

\input{tables/table_os}

\section{Benchmark Tasks and Experimental Setup}
The three tasks (Figure~\ref{fig:tasks}) each isolate one property of real logging. All use accuracy-style recall at rank $K$ (R@$K$, $K\in\{1,3,5\}$) over the 218 classes, and F1 and precision at rank one (P@1) where indicated; we report the scores as produced by single evaluations without confidence intervals. The material available to us does not specify the train/test split, the F1 averaging scheme, the unit on which Task-2 accuracy is computed, the window and step size of the temporal protocol, how User~A was selected, or the fine-tuning hyperparameters; we do not invent them here, and reproducing Task~2 requires them.

\begin{figure}[!tb]
  \centering
  \includegraphics[width=\columnwidth]{ACM_task.jpg}
  \caption{The three evaluation tasks: closed-set classification of cropped dishes by embedding similarity (Task 1), continual user-aware adaptation along the timeline (Task 2) and menu-aware classification of a dish within a multi-dish image (Task 3).}
  \label{fig:tasks}
\end{figure}

\textbf{Task 1: closed-set classification.} The input is a dish crop given by its bounding box and the output is one of the 218 labels. Contrastive models (CLIP ViT-L/14~\cite{CLIP}, EVA-CLIP ViT-g~\cite{sun2023eva}) rank labels by image--text similarity. For generative LMMs, which do not output a ranking over a closed label set, a free-form dish description is generated, embedded with intfloat/e5-large-v2~\cite{wang2023e5}, and the label whose embedding has the highest cosine similarity is selected, with the ranking used for R@3 and R@5. The LMMs are evaluated through lmms-eval~\cite{lmms_eval2024}: GPT-4o (a GPT-4-family model~\cite{achiam2023gpt}), Qwen2.5-VL-7B-Instruct, Qwen2-VL-7B-Instruct, Qwen2-VL-3B-Instruct, llava-onevision-qwen2-7b-ov, llava-v1.6-34b, llava-v1.5-7b~\cite{liu2024improved} and llava-v1.6-mistral-7b.

\textbf{Task 2: continual user-aware adaptation.} Records are ordered by date. A model is trained on the data up to a point in time and tested on the following data, and the training set grows over time (Figure~\ref{fig:tasks}). Accuracy is reported per week together with a one-month moving average ($\pm2$ weeks) weighted in proportion to the weekly data volume. We compare a zero-shot model with full fine-tuning and LoRA fine-tuning~\cite{hu2022lora} of Qwen2.5-VL-3B on all users, and we study personalization by additionally fine-tuning a LoRA on a single heavy logger (User~A).

\textbf{Task 3: menu-aware classification.} In multi-dish images, a dish may be easier to identify given the rest of the meal. A model first produces a description of the entire image (the menu), and then classifies the cropped dish with the description as context (Figure~\ref{fig:desc}), using Qwen2.5-VL-7B-Instruct. We report F1, P@1 and R@1/3/5 with and without the description.

\begin{figure}[!tb]
  \centering
  \includegraphics[width=\columnwidth]{description_eg.jpg}
  \caption{Description-guided classification. Left: the description fixes an error made on the crop alone. Right: the description introduces an error on a crop that was correct without it.}
  \label{fig:desc}
\end{figure}

\section{Benchmark Validation}
\textbf{Difficulty relative to a web dataset.} To check that realistic data are harder than a curated web dataset, we compare feature-space statistics of FoodLogAthl-218 and UEC-Food 256~\cite{kawano2015automatic} with three feature extractors (CLIP~\cite{CLIP}, EVA-CLIP~\cite{sun2023eva} and Swin Transformer~\cite{liu2021swin}): the trace of the feature spread, the Silhouette score~\cite{rousseeuw1987silhouettes}, the Davies--Bouldin index~\cite{davies1979cluster} and the Calinski--Harabasz index~\cite{calinski1974dendrite}. Table~\ref{table:diversity_metrics_comparison} shows that FoodLogAthl-218 has a larger spread and less compact, more overlapping classes for all four statistics and all three extractors. For instance, the trace with CLIP is 81.70 against 71.21, and the Silhouette score is $-0.021$ against $0.033$, i.e., classes overlap on average. These statistics describe the geometry of the classes and not model accuracy, and they compare two datasets that differ in size and in number of classes (218 against 256), so they support greater intra-class diversity and class overlap but are not a calibrated difficulty measure.

\input{tables/table_xj}

\textbf{Annotation quality.} Labels come from users and were then checked by automated outlier detection, manual review of boxes and labels, and human adjudication of the GPT-4o label groups. The manual review removed 612 images and 3,164 boxes (step 5 of Table~\ref{tab:filtering}). We did not measure inter-annotator agreement, so the residual label noise is unquantified. Label merging reduces the label space from 1,180 to 218 but makes some classes coarse, and some dishes may remain ambiguous between classes.

\textbf{Sanity checks.} The benchmark is not trivial (best top-1 recall 42.4\%, Table~\ref{table:classification_comparison}) and not degenerate (models differ strongly and rankings depend on the metric). We have not run a leakage check against the pre-training data of the evaluated models, which cannot be excluded for web-scale models, though the images come from private logs and not from the web.

\section{Benchmark Results}
\subsection{Task 1: Closed-Set Classification}
Table~\ref{table:classification_comparison} shows that classification of real logged dishes is hard for all ten systems. GPT-4o is the best on F1 (0.373), R@1 (0.424) and R@3 (0.499) but its R@5 (0.522) is below that of both contrastive models; CLIP ViT-L/14 has the best R@5 (0.600), followed by EVA-CLIP (0.594). The ranking therefore changes with the metric. At rank one, four LMMs (GPT-4o, Qwen2.5-VL-7B, Qwen2-VL-7B and llava-onevision) exceed the contrastive models (0.363--0.424 against 0.257--0.265), whereas the contrastive models are ahead of all LMMs except GPT-4o at R@3 and ahead of all LMMs at R@5 (0.594--0.600 against at most 0.522), and their F1 (0.239--0.247) is higher than that of every LMM except GPT-4o. Within a family, scale matters: Qwen2-VL-7B improves over Qwen2-VL-3B by about 10 points at R@1, and the three LLaVA-1.5/1.6 variants have the lowest R@1 (0.193--0.243). The gap between R@1 and F1 is large for several LMMs (e.g., Qwen2-VL-7B, 0.366 against 0.159), which suggests that some classes are predicted poorly even when top-1 recall is moderate, as expected for a long-tailed label set; we did not analyze per-class errors, so this explanation is a hypothesis.

\input{tables/table_kf}

\subsection{Task 2: Continual User-Aware Adaptation}
\begin{figure*}[!tb]
  \centering
  \begin{minipage}[b]{0.49\textwidth}\centering
    \includegraphics[width=\linewidth]{all_data_1_35_with_legend.jpg}\\[-1pt]
    {\small (a) All users}
  \end{minipage}\hfill
  \begin{minipage}[b]{0.49\textwidth}\centering
    \includegraphics[width=\linewidth]{user_zero_lora_ma_wide_legend.jpg}\\[-1pt]
    {\small (b) User~A}
  \end{minipage}
  \caption{Task 2 accuracy over time. (a) Zero-shot, full fine-tuning and LoRA fine-tuning on all users. (b) User~A: zero-shot, LoRA fine-tuned on all users, and LoRA fine-tuned on User~A only. Dashed lines are weekly values and solid lines the weighted moving average. Tick labels are as in the source figures (in (b) the label 2023/10 appears twice), so month positions are approximate.}
  \label{fig:task2}
\end{figure*}

Figure~\ref{fig:task2}a shows weekly accuracy and its weighted moving average for Qwen2.5-VL-3B. With the moving average, the zero-shot model starts near 0.23 and declines to about 0.16 at the end of the plotted period, whereas full fine-tuning and LoRA fine-tuning increase from about 0.25 to about 0.5 (values read from the figure). Full fine-tuning is above LoRA for most of the period and the two are close in the last weeks. At the end the fine-tuned models are about 0.3 above the zero-shot model. Weekly values (dashed) fluctuate by more than 0.1, so differences between full and LoRA fine-tuning in individual weeks should not be interpreted. Why the zero-shot model gets worse is not established by our data; a change in the composition of the weekly data is a possibility that we have not tested.

For one heavy logger (Figure~\ref{fig:task2}b), a LoRA trained only on that user's data goes from about 0.17 in the first weeks to about 0.6--0.67 later, a gain of roughly 45--50 points (and about 0.5 above the zero-shot model), and is above the all-user LoRA from about July to December 2023, while the zero-shot model remains between about 0.05 and 0.14. Early in the stream, however, the all-user adapter is higher (about 0.3 against 0.2), and at the last plotted point it is higher again (about 0.65 against 0.57), so the personalized adapter is not uniformly better. This is the result for a single user, and the weekly values (dashed) jump between 0 and 0.8, so the curves are noisy; it should be read as a demonstration of what the protocol can measure and not as an estimate of the benefit of personalization in general.

\subsection{Task 3: Menu-Aware Classification}
Table~\ref{table:desc_guided} shows that adding the whole-image description raises F1 from 0.159 to 0.230 and P@1 from 0.102 to 0.179, but lowers R@1 from 0.366 to 0.324, R@3 from 0.459 to 0.403 and R@5 from 0.478 to 0.434. At the level of individual predictions the description corrects 687 errors but also causes 1,222 new errors, a net increase of 535 errors. Context is thus neither uniformly helpful nor harmful: it changes the prediction behaviour in a way that raises precision but loses recall (Figure~\ref{fig:desc}). The ``w/o description'' row coincides numerically with the Qwen2-VL-7B-Instruct row, not the Qwen2.5-VL-7B-Instruct row (F1 0.235), of Table~\ref{table:classification_comparison}; we report the numbers as supplied, treat the comparison as internal to Task 3, and make no claim of improvement over the Task-1 result (0.2304 is 0.0045 below 0.2349).

\input{tables/table_vs}

\section{Discussion}
\textbf{What the benchmark offers.} FoodLogAthl-218 measures food recognition under conditions of real logging: natural imbalance, casual photographs, temporal order and multi-dish meals. The three tasks reveal what a single closed-set accuracy would hide: rankings differ between R@1 and R@5 and between F1 and recall, adaptation along time matters strongly for a fixed model, and extra context can trade recall for precision. We therefore recommend reporting several metrics together.

\textbf{What it does not show.} The data come from one application and 513 users after filtering, with labels chosen from one database; the geographic and cultural distribution of the users is not characterised here, so conclusions about other populations are not warranted. The curation removed 40\% of images and all classes with fewer than 10 samples, so very rare dishes are out of scope, and label merging makes the label set coarser than the 140K-item database. Inter-annotator agreement was not measured and the residual noise is unknown. The Task-2 results use one model family and one heavy user and are read from plots without confidence intervals. The baselines are off-the-shelf models queried with a single pipeline, so prompts or the text-embedding mapping might alter the ranking. Finally, accuracy here says nothing about calorie or nutrition estimation, the downstream goal of logging. The dataset is intended for research on closed-set, continual, personalized and context-aware recognition, under the privacy protections (two-stage removal of personal information, informed consent) with which the data were collected.

\section{Conclusion}
We introduced FoodLogAthl-218, a benchmark of 6,925 images, 14,349 dish bounding boxes and 218 classes built from real dietary logs, together with three tasks and a curation pipeline whose steps are documented. Ten off-the-shelf models reach at most 42.4\% R@1 and 60.0\% R@5 on closed-set classification, the data are more diverse and overlapping than UEC-Food 256 in feature space, fine-tuning along the timeline lifts accuracy by roughly 0.3 over zero-shot, a user-specific adapter helps one heavy logger for much of the period, and menu context trades recall for precision. These results show that realistic logging data expose gaps that web-collected benchmarks do not, while the benchmark's single-source origin and unmeasured annotation agreement bound the conclusions that can be drawn from it.

\bibliographystyle{ACM-Reference-Format}

\bibliography{references}

\end{document}
